\chapter{Jesus subiu aos céus, está sentado à direita de Deus Pai Todo-Poderoso}
\chaptermark{Jesus subiu aos céus}

O sexto artigo da Profissão de Fé proclama a glorificação definitiva da humanidade de Cristo: ``Jesus subiu aos céus, está sentado à direita de Deus Pai Todo-Poderoso''. Este mistério completa o Mistério Pascal, marcando a entrada irreversível da humanidade de Jesus na glória divina. A Ascensão não representa o fim da presença de Cristo, mas o início de um novo modo de presença. O Catecismo da Igreja Católica, nos parágrafos 659 a 667, apresenta esta verdade como o coroamento da obra da salvação, explicando o seu significado teológico e as suas consequências para a vida da Igreja e de cada fiel.

\section{Jesus subiu aos céus (659--662)}

A Ascensão do Senhor é o acontecimento que encerra as aparições do Ressuscitado e inaugura a sua glorificação plena. O corpo de Cristo foi glorificado desde o momento da Ressurreição, como demonstram as propriedades novas e sobrenaturais que passou a possuir. No entanto, durante os quarenta dias, a sua glória permaneceu velada sob a aparência de humanidade ordinária, enquanto comia e bebia com os discípulos e os instruía sobre o Reino de Deus.

A última aparição termina com a entrada irreversível da sua humanidade na glória divina, simbolizada pela nuvem e pelo céu. Jesus foi elevado à vista dos discípulos e uma nuvem o subtraiu aos seus olhos. Este acto não foi obra de uma força exterior, mas o culminar da sua própria vontade e do poder divino. Cristo, que desceu do céu na Encarnação, regressa agora ao Pai com a sua humanidade glorificada.

A Ascensão está intimamente ligada à Encarnação. Somente Aquele que veio do Pai podia regressar ao Pai. Ninguém pode subir ao céu por suas próprias forças. Apenas Cristo, que desceu do céu, pode abrir-nos o acesso à casa do Pai. A sua Ascensão é, portanto, garantia da nossa própria ascensão futura. O Prefácio da Ascensão afirma com clareza: ``para que, onde a Cabeça nos precedeu, nós, seus membros, tenhamos a firme esperança de a seguir''.

Pela Ascensão, Jesus entra no santuário celeste, não feito por mãos humanas, para interceder por nós permanentemente. Ele é o Sumo Sacerdote que se apresenta agora perante o Pai em nosso favor. A sua humanidade glorificada tornou-se o centro da liturgia celestial, onde Ele continua a exercer o seu sacerdócio eterno.

A Ascensão revela também o sentido da Cruz. Quando Jesus afirma ``e eu, quando for elevado da terra, atrairei todos a mim'', refere-se simultaneamente à Cruz e à Ascensão. O seu levantamento na Cruz inicia o seu levantamento para a glória do céu.

\section{Está sentado à direita de Deus Pai Todo-Poderoso (663--667)}

A expressão ``sentado à direita do Pai'' não indica uma posição espacial, mas a glória e a honra da divindade. Aquele que é Filho de Deus desde toda a eternidade, depois de se ter feito homem e de ter glorificado a sua carne, senta-se corporalmente à direita do Pai. Esta expressão bíblica significa que Jesus partilha plenamente o poder, a honra e a glória do Pai.

Estar sentado à direita do Pai significa o início do Reino messiânico. Cumpre-se a visão de Daniel sobre o Filho do Homem a quem foi dado domínio, glória e reino. O Reino de Cristo não terá fim. Os Apóstolos tornaram-se testemunhas deste Reino que durará eternamente.

A Ascensão e a sessão à direita do Pai marcam a inauguração da realeza de Cristo. Ele reina agora sobre toda a criação. Todos os poderes, autoridades e domínios lhe estão submetidos. Ele é a Cabeça da Igreja, que é o seu Corpo e a plenitude daquele que tudo plenifica.

Da sua posição à direita do Pai, Cristo intercede continuamente por nós. Ele vive para sempre para interceder em nosso favor. Como Sumo Sacerdote dos bens futuros, Ele é o centro da liturgia que honra o Pai no céu. A sua intercessão permanente garante a efusão contínua do Espírito Santo sobre a Igreja.

A Ascensão não afasta Cristo da Igreja. Pelo contrário, torna-o presente de modo novo e mais profundo. Embora escondido aos olhos dos homens, Ele permanece presente na sua Igreja, especialmente nos sacramentos. A nuvem que o ocultou aos discípulos é o mesmo véu que nos permite viver pela fé e não pelo que se vê.

A fé na Ascensão e na sessão à direita do Pai enche o cristão de esperança. Cristo, nossa Cabeça, precedeu-nos na glória do Pai para que nós, membros do seu Corpo, tenhamos a firme confiança de o seguir um dia. Ele não nos deixou órfãos. Da sua glória, envia o Espírito Santo e prepara-nos um lugar na casa do Pai.

Este mistério convida-nos a orientar o nosso olhar para o alto, sem deixar de trabalhar com dedicação na terra. A Ascensão não é fuga do mundo, mas a garantia de que o nosso destino final está unido a Cristo glorioso. Enquanto esperamos a sua manifestação final, vivemos na esperança e no testemunho alegre da fé pascal.